\chapter{Image Manipulation and Text}

\section{NLP Basics}
\dfn{Natural Language Processing}{
Subfield of AI/linguistics/computer science focused on processing and analyzing natural language data.
}
\dfn{Tokenization}{
Split text into minimal units (words/subwords/sentences) for downstream processing.
}
\dfn{Stemming vs Lemmatization}{
Stemming chops affixes to crude roots; lemmatization maps to dictionary form using vocab and POS.
}
\dfn{POS Tagging}{
Assign grammatical category (noun, verb, adj, etc.) to each token.
}
\dfn{Named Entity Recognition}{
Locate and classify spans as persons, locations, orgs, dates, etc.
}
\dfn{Sentiment Analysis}{
Classify polarity of text (positive/negative/neutral).
}
\dfn{Language Modeling}{
Estimate $p(w_t \mid w_{<t})$ to predict next token (autocomplete, MT).
}
\dfn{Machine Translation}{
Map sequence in source language to target language automatically.
}
\dfn{Text Summarization}{
Produce a concise version preserving key information (extractive or abstractive).
}
\dfn{Information Extraction}{
Pull structured facts (entities/relations/events) from unstructured text.
}
\dfn{Text Classification}{
Assign documents to predefined labels (spam, topic, sentiment).
}
\dfn{Topic Modeling}{
Unsupervised discovery of latent themes (e.g., LDA).
}
\dfn{Dependency Parsing}{
Recover syntactic head/dependent relations between words.
}
\dfn{Discourse Analysis}{
Study coherence across sentences (rhetorical structure, cohesion).
}
\dfn{Regular Expressions}{
Patterns for searching/extracting/editing text programmatically.
}
\dfn{Web Scraping}{
Programmatic retrieval/parsing of web pages (e.g., BeautifulSoup over HTML).
}

\section{NLP Code Snippets (Python)}
\subsection*{Tokenize, Stem, Lemmatize (NLTK)}
\begin{lstlisting}
import nltk
from nltk.stem import PorterStemmer, WordNetLemmatizer
nltk.download('punkt'); nltk.download('wordnet')

text = "Cats are happily running."
tokens = nltk.word_tokenize(text)
stemmer = PorterStemmer()
lemmatizer = WordNetLemmatizer()

stems = [stemmer.stem(t) for t in tokens]
lems  = [lemmatizer.lemmatize(t, pos="v") for t in tokens]
print(tokens); print(stems); print(lems)
\end{lstlisting}

\subsection*{POS \& NER (spaCy)}
\begin{lstlisting}
import spacy
nlp = spacy.load("en_core_web_sm")
doc = nlp("Apple is looking at buying U.K. startup for $1B.")
for tok in doc:
    print(tok.text, tok.pos_, tok.dep_, tok.head.text)
for ent in doc.ents:
    print(ent.text, ent.label_)
\end{lstlisting}

\subsection*{Regex extraction}
\begin{lstlisting}
import re
text = "Contact us at support@example.com or hr@example.org"
emails = re.findall(r"[\\w.-]+@[\\w.-]+", text)
print(emails)
\end{lstlisting}

\subsection*{Bag-of-Words + TF-IDF + Linear Classifier}
\begin{lstlisting}
from sklearn.feature_extraction.text import TfidfVectorizer
from sklearn.linear_model import LogisticRegression
from sklearn.pipeline import make_pipeline

X = ["I love this movie", "Terrible acting", "Great plot", "Worst film ever"]
y = [1, 0, 1, 0]  # 1=pos, 0=neg
clf = make_pipeline(TfidfVectorizer(), LogisticRegression())
clf.fit(X, y)
print(clf.predict(["fantastic performance", "awful script"]))
\end{lstlisting}

\section{Computer Vision with OpenCV}
\dfn{OpenCV}{
Open-source CV library for image I/O, transformations, and basic vision ops.
}

\subsection*{Quick Notes from Class Demos}
\begin{itemize}
    \item Images are NumPy arrays in BGR order; \texttt{img.shape} gives (rows, cols, channels).
    \item Display with \texttt{cv2.imshow} + \texttt{waitKey} + \texttt{destroyAllWindows}; Matplotlib expects RGB.
    \item Channels: \texttt{b,g,r = cv2.split(img)} then \texttt{cv2.merge([...])}.
    \item Resize: \texttt{cv2.resize(img,(w,h))} or scale factors \texttt{fx, fy}; \texttt{INTER\_CUBIC} for up, \texttt{INTER\_AREA} for down.
    \item Colorspace: \texttt{cv2.cvtColor(img, cv2.COLOR\_BGR2GRAY/RGB)}.
    \item Arithmetic: \texttt{cv2.add} or \texttt{cv2.addWeighted(img1,a,img2,b,gamma)}.
    \item Filters: \texttt{cv2.blur}, \texttt{cv2.GaussianBlur}, \texttt{cv2.medianBlur}.
    \item Rotation: \texttt{M=cv2.getRotationMatrix2D(center, angle, scale)} then \texttt{cv2.warpAffine}.
    \item ROI crop: slice \texttt{img[y1:y2, x1:x2]}.
    \item Load from URL: read bytes, \texttt{np.asarray}, then \texttt{cv2.imdecode}.
\end{itemize}

\subsection*{Read, Show, Save}
\begin{lstlisting}
import cv2
img = cv2.imread("input.jpg")              # BGR
cv2.imshow("img", img); cv2.waitKey(0)
cv2.imwrite("out.png", img)
\end{lstlisting}

\subsection*{Resize}
\begin{lstlisting}
resized = cv2.resize(img, (256, 256), interpolation=cv2.INTER_AREA)
\end{lstlisting}

\subsection*{Scale Up / Down}
\begin{lstlisting}
# upsample 2x
up = cv2.resize(img, None, fx=2, fy=2, interpolation=cv2.INTER_CUBIC)
# downsample 0.5x
down = cv2.resize(img, None, fx=0.5, fy=0.5, interpolation=cv2.INTER_AREA)
\end{lstlisting}

\subsection*{Split / Merge Channels}
\begin{lstlisting}
b, g, r = cv2.split(img)
merged = cv2.merge([b, g, r])
\end{lstlisting}

\subsection*{Display Each Channel}
\begin{lstlisting}
cv2.imshow("b", b); cv2.waitKey(0)
cv2.imshow("g", g); cv2.waitKey(0)
cv2.imshow("r", r); cv2.waitKey(0); cv2.destroyAllWindows()
\end{lstlisting}

\subsection*{Blend / Add Images}
\begin{lstlisting}
img2 = cv2.imread("overlay.png")
blend = cv2.addWeighted(img, 0.7, img2, 0.3, 0)
\end{lstlisting}

\subsection*{Rotate}
\begin{lstlisting}
h, w = img.shape[:2]
M = cv2.getRotationMatrix2D((w/2, h/2), 45, 1.0)
rotated = cv2.warpAffine(img, M, (w, h))
\end{lstlisting}

\subsection*{ROI (Crop)}
\begin{lstlisting}
roi = img[100:200, 150:300]  # y1:y2, x1:x2
\end{lstlisting}

\subsection*{Color Space Conversion}
\begin{lstlisting}
gray = cv2.cvtColor(img, cv2.COLOR_BGR2GRAY)
\end{lstlisting}

\subsection*{Filtering}
\begin{lstlisting}
avg    = cv2.blur(img, (3,3))
gauss  = cv2.GaussianBlur(img, (3,3), 0)
median = cv2.medianBlur(img, 3)
\end{lstlisting}

\subsection*{Load from URL and Crop ROI}
\begin{lstlisting}
import urllib.request as req, numpy as np
data = req.urlopen("https://upload.wikimedia.org/wikipedia/commons/7/78/Pisa_-_veduta_aerea_2.JPG").read()
arr = np.asarray(bytearray(data), dtype=np.uint8)
remote = cv2.imdecode(arr, -1)
roi = remote[100:350, 1200:1800]
\end{lstlisting}

\section{Practice / Exam Prompts}
\pr{When to stem vs lemmatize}{
Stemming is faster/rougher (search recall). Lemmatization is slower but produces valid lemmas (better precision/NLU).
}
\pr{Regex task}{
Write a pattern to extract dates like 2025-04-11 or 11/04/2025: e.g., `(\\d{4}-\\d{2}-\\d{2})|(\\d{2}/\\d{2}/\\d{4})`.
}
\pr{Design a sentiment pipeline}{
Steps: tokenize $\rightarrow$ lowercase/clean $\rightarrow$ vectorize (TF-IDF) $\rightarrow$ train logistic regression; evaluate with held-out set.
}
\pr{Image resize + rotate}{
Given a large photo, resize to 512x512, then rotate 30° about center using OpenCV (`cv2.resize`, `getRotationMatrix2D`, `warpAffine`).
}
\pr{ROI use case}{
Crop a face region for downstream model: use slicing or `cv2.selectROI` to extract and save the patch.
}
